\section{Design decision log}
\label{sec:decisions}

Append-only log. Newest entries at the bottom. If a decision is reversed, add a
new entry that supersedes the old id; do not silently delete history.

%----------------------------------------------------------------------------
\subsection*{D-001 --- Package rename (structure)}
\begin{decisionbox}[title={D-001 \quad 2026-07-28}]
\textbf{Decision:} Rename/split ambiguous folders to match responsibility:
\texttt{Core$\rightarrow$Game}, \texttt{System$\rightarrow$\{App,Engine,Services\}},
\texttt{Init+Util$\rightarrow$Foundation}, \texttt{AssetLoaders$\rightarrow$Assets}.\\
\textbf{Why:} Names were lying; slowed reasoning and migrations.\\
\textbf{Consequences:} Include/CMake paths updated; archive for old States.
\end{decisionbox}

%----------------------------------------------------------------------------
\subsection*{D-002 --- MathTypes header name}
\begin{decisionbox}[title={D-002 \quad 2026-07-28}]
\textbf{Decision:} Project math aliases live in \pathref{MathTypes.h}, not
\pathref{Math.h}.\\
\textbf{Why:} Windows case-insensitive include collision with CRT
\texttt{math.h}.\\
\textbf{Consequences:} All includes updated; the Foundation package map under
\pathref{docs/packages/} notes the trap.
\end{decisionbox}

%----------------------------------------------------------------------------
\subsection*{D-003 --- Render abstraction direction}
\begin{decisionbox}[title={D-003 \quad 2026-07-28 (implemented; see D-R*) }]
\textbf{Decision:} Command/token submission via \symbolref{IBackend}
+ \symbolref{RenderCommand} queue; game code must not permanently depend on
raw GL for draw submission.\\
\textbf{Why:} Multi-API optionality; clearer frame boundary; testability.\\
\textbf{Consequences:} Phases 1--6 complete for planned scope
(\cref{sec:migration-phases}); hybrid immediate \texttt{Draw} remains only as
fallback.
\end{decisionbox}

%----------------------------------------------------------------------------
\subsection*{D-004 --- Module frame hooks}
\begin{decisionbox}[title={D-004 \quad inferred from IModule}]
\textbf{Decision:} Modules expose \symbolref{Start}/\symbolref{Update}/
\symbolref{DispatchDrawables}/\symbolref{Exit}.\\
\textbf{Why:} Separate simulation tick from draw contribution.\\
\textbf{Consequences:} Draw list is rebuilt or contributed each frame through
dispatch. Module ownership/composition was later closed by D-E1.
\end{decisionbox}

%----------------------------------------------------------------------------
\subsection*{D-005 --- Rules as polymorphic sets}
\begin{decisionbox}[title={D-005 \quad inferred}]
\textbf{Decision:} Each CA variant is a \symbolref{RuleSet} subclass.\\
\textbf{Why:} Open for extension; shared generation sweep.\\
\textbf{Consequences:} Factory logic currently in \symbolref{CellContext}
string matching (could become a registry later).
\end{decisionbox}

%----------------------------------------------------------------------------
\subsection*{D-006 --- Mode handling (enum vs State classes)}
\begin{decisionbox}[title={D-006 \quad current code; author may supersede}]
\textbf{Decision (current):} \symbolref{CellState} enum inside
\symbolref{CellGameModule}.\\
\textbf{Supersedes:} class hierarchy under \pathref{archive/old-states/}.\\
\textbf{Why (inferred):} fewer types for a small mode set.\\
\textbf{Open at decision time:} SAVE/LOAD completeness; whether splash/global
state is acceptable.\\
\textbf{Resolution (2026-08-04):} frame modes are now
\texttt{NORMAL | EDIT | EXIT}; save/load are registered commands with validated
file handling (D-CLI1).
\end{decisionbox}

%----------------------------------------------------------------------------
\subsection*{D-007 --- Enrollment vs per-frame commands}
\begin{decisionbox}[title={D-007 \quad visible in Renderer comments}]
\textbf{Decision:} Resource creation goes through enroll/Create* APIs; do not
stuff mesh/shader creation into the per-frame token stream.\\
\textbf{Why:} Keep the frame queue about \emph{binding and drawing}.\\
\textbf{Consequences:} Need an explicit destroy/reload story for hot-reload
later.
\end{decisionbox}

%----------------------------------------------------------------------------
\subsection*{D-008 --- contributing constraints}
\begin{decisionbox}[title={D-008 \quad from docs/contributing.md}]
\textbf{Decision:} Avoid \texttt{auto}; avoid namespaces (prefer static
classes/structs); no recursion; third-party via PR.\\
\textbf{Why:} Author house style / learning constraints.\\
\textbf{Consequences:} Documentation and new code should respect these unless
explicitly waived in a later decision.
\end{decisionbox}

%----------------------------------------------------------------------------
\subsection*{D-R1 --- RenderCommand payload shape}
\begin{decisionbox}[title={D-R1 \quad 2026-07-28}]
\textbf{Decision:} \symbolref{RenderCommand} uses a \textbf{simple tagged union}
(\texttt{CommandType} + POD payload union), not a kitchen-sink flat struct and
not a byte-stream compiler. Prefer C-style union over \texttt{std::variant}.\\
\textbf{Why:} Type-safe uniforms/clear/update data while remaining
copyable for \symbolref{ArrayQueue}; matches house style.\\
\textbf{Consequences:} \pathref{Rendering/RenderCommand.h} redesigned; execute
path switches on type and reads the matching payload.
\end{decisionbox}

%----------------------------------------------------------------------------
\subsection*{D-R2 --- Who emits tokens}
\begin{decisionbox}[title={D-R2 \quad 2026-07-28}]
\textbf{Decision:} Each drawable emits its own tokens via
\texttt{AppendCommands(Renderer*)} (or equivalent). \symbolref{Renderer}
owns frame setup (clear/viewport) and submit.\\
\textbf{Why:} Smallest strangler step from today's
\symbolref{DispatchDrawables} list.\\
\textbf{Consequences:} Hybrid phase may still call immediate \texttt{Draw}
for unmigrated drawables after token submit.
\end{decisionbox}

%----------------------------------------------------------------------------
\subsection*{D-R3 --- Handle typing for v1}
\begin{decisionbox}[title={D-R3 \quad 2026-07-28}]
\textbf{Decision:} Keep \texttt{unsigned long} opaque handles with backend
registries (\texttt{tableID} $\rightarrow$ GL object). Strong typedefs deferred.\\
\textbf{Why:} Minimal churn; critical fix is resolve-on-submit, not type names.\\
\textbf{Consequences:} Execute must never treat handles as raw GLuints without
registry lookup.
\end{decisionbox}

%----------------------------------------------------------------------------
\subsection*{D-R4 --- Migration scope}
\begin{decisionbox}[title={D-R4 \quad 2026-07-28}]
\textbf{Decision:} Full migration Phases 1--5 (token proof $\rightarrow$ Scene
routing $\rightarrow$ Canvas $\rightarrow$ UI drawables $\rightarrow$ cleanup).
Phase 6 (mock backend) included; second real GPU API still optional.\\
\textbf{Why:} Playable strangler path to API isolation without big-bang rewrite.\\
\textbf{Also freeze for v1:} destroy resources at shutdown only; Windows GL is
primary platform.\\
\textbf{Superseded in part by D-P1/D-P4:} full-frame Canvas rewrite was v1 only.
Dirty-rect/PBO uploads remain, while the R8 presentation was later superseded
by the restored RGB fade path documented in D-P4.
\end{decisionbox}

%----------------------------------------------------------------------------
\subsection*{D-R5 --- Blend state on first enable}
\begin{decisionbox}[title={D-R5 \quad 2026-07-28}]
\textbf{Decision:} When enabling blending in \symbolref{GLDevice::ApplyPipelineState},
always call \texttt{glBlendFunc} (do not rely on CPU-side factor cache matching
defaults).\\
\textbf{Why:} OpenGL default is \texttt{ONE}/\texttt{ZERO}; cache defaults were
SrcAlpha/OneMinusSrcAlpha so factors looked ``unchanged'' and console panel
lost transparency after the token migration.\\
\textbf{Consequences:} Semi-transparent UI (console panel alpha) works again.
\end{decisionbox}

%----------------------------------------------------------------------------
\subsection*{D-R6 --- Dead stack cleanup (Phase 5)}
\begin{decisionbox}[title={D-R6 \quad 2026-07-28}]
\textbf{Decision:} Archive window \symbolref{RenderQueue}, legacy
\texttt{Drawable.cpp} GL helpers, empty \texttt{Transform}/\texttt{RenderContext},
and unused \texttt{CellCommandLine} under \pathref{archive/dead-render/}.
\symbolref{DrawableBase} is a pure list/token interface (no GL handles).\\
\textbf{Why:} Token path won; dual submission models confuse maintenance.\\
\textbf{Consequences:} \texttt{IRenderWindow::getRenderQueue} removed; frame
draw list is only \symbolref{Scene} + \symbolref{Renderer}.
\end{decisionbox}

%----------------------------------------------------------------------------
\subsection*{D-R7 --- Mock backend for tests (Phase 6)}
\begin{decisionbox}[title={D-R7 \quad 2026-07-28}]
\textbf{Decision:} Add headless \texttt{MockBackend} implementing
\symbolref{IBackend}: record \texttt{Create*} and snapshot the command queue
on \texttt{SubmitCommandQueue}. Originally shipped as CMake target
\texttt{IllumoRenderTests} / CTest \texttt{MockBackend}; the current canonical
target and CTest name are \texttt{IllumoTests}, with the old names retained only
as compatibility aliases. Do \textbf{not} add Vulkan/Metal yet.\\
\textbf{Why:} First meaningful automated render test without a GPU context;
proves token vocabulary and handle plumbing independently of GL.\\
\textbf{Consequences:} CI can run \texttt{ctest -R IllumoTests}; second real API
remains future work after more production needs appear.
\end{decisionbox}

%----------------------------------------------------------------------------
\subsection*{D-R8 --- Inject IBackend into Renderer}
\begin{decisionbox}[title={D-R8 \quad 2026-07-28, refined 2026-08-07}]
\textbf{Decision:} \symbolref{Renderer} is backend-neutral: its only constructor
takes an injected \symbolref{IBackend*} and a \texttt{takeOwnership} flag.
Production composition (\symbolref{Illumo}\texttt{::Init}) constructs the real
backend via \texttt{CreateOpenGLBackend} under \pathref{Rendering/OpenGL/} and
injects it with \texttt{takeOwnership=true}. Tests inject \texttt{MockBackend}
with \texttt{takeOwnership=false}. \texttt{Renderer.h} must not include or
construct concrete OpenGL types.\\
\textbf{Why:} End-to-end drawable $\rightarrow$ enroll/token $\rightarrow$ submit
without a GPU; second platforms can supply a different \symbolref{IBackend}
without touching the generic frame-submit surface.\\
\textbf{Consequences:} \texttt{TestRendererE2E} covers inject, owned-backend
composition style, header neutrality, \texttt{RenderScene}, Canvas tokens,
hybrid immediate drawables, and \texttt{RenderProofQuad}. OpenGL sources stay
under \pathref{Rendering/OpenGL/}; Mock under \pathref{Rendering/Mock/}.
\end{decisionbox}

%----------------------------------------------------------------------------
\subsection*{D-P1 --- Dirty visual path}
\begin{decisionbox}[title={D-P1 \quad 2026-07-28}]
\textbf{Decision:} Canvas/CellGameModule use dirty flags so idle frames skip
full-grid visual work and GPU upload.\\
\textbf{Why:} Dominant cost was recoloring + full texture upload every render
frame even when the CA did not step.\\
\textbf{Consequences:} \texttt{cellsDirty}, upload pending, early-out
\texttt{tickVisual}; still life after sim step does not re-dirty.
\end{decisionbox}

%----------------------------------------------------------------------------
\subsection*{D-P2 --- UI batching and geometry cache}
\begin{decisionbox}[title={D-P2 \quad 2026-07-28}]
\textbf{Decision:} CommandLine emits one packed \texttt{UpdateBuffer}/draw when
open; GLString caches GPU geometry until content/style changes.\\
\textbf{Why:} Per-line uploads and FPS string rebuilds were measurable when the
console/FPS overlay was active.\\
\textbf{Consequences:} Token stream stays small; FPS may rebuild $\sim$1\,Hz.
\end{decisionbox}

%----------------------------------------------------------------------------
\subsection*{D-P3 --- Double-buffer CA generation}
\begin{decisionbox}[title={D-P3 \quad 2026-07-28}]
\textbf{Decision:} \texttt{RuleSet::calcGeneration} writes pure
\texttt{nextState(cell, aliveNeighbors)} into a scratch buffer, then
\texttt{memcpy}s back; neighbor counts use raw grid pointers + toroidal wrap.\\
\textbf{Why:} Correct simultaneous update; faster than per-cell
\texttt{getCanvasPixel}/in-place virtual writes; sparse dirty AABB for visuals.\\
\textbf{Consequences:} All active rulesets override \texttt{nextState}; stubs
use identity.\\
\textbf{Refined by D-P5 (2026-08-06):} scratch lives on \symbolref{CellGrid}
as a back buffer; promotion is pointer swap, not full-grid \texttt{memcpy}.
\end{decisionbox}

%----------------------------------------------------------------------------
\subsection*{D-P5 --- Single-pass dirty AABB + buffer swap}
\begin{decisionbox}[title={D-P5 \quad 2026-08-06}]
\textbf{Decision:} Track changed-cell AABB during the eval pass; store next
generation in \texttt{CellGrid} back buffer; on any change call
\texttt{swapLifeBuffers()} and mark the sparse dirty region. Interior Moore
counts skip toroidal wrap branches.\\
\textbf{Why:} Second full-grid compare pass and full \texttt{memcpy} write-back
dominated serial generation cost on large dense grids.\\
\textbf{Consequences:} \texttt{lifeCanvas} always means current front; rules
must not retain raw pointers across a swap; still life leaves domain clean.
\end{decisionbox}

%----------------------------------------------------------------------------
\subsection*{D-P6 --- Fade loop + packed dirty PBO staging}
\begin{decisionbox}[title={D-P6 \quad 2026-08-06}]
\textbf{Decision:} Keep CPU RGB fade presentation. Hoist coordinates in
\texttt{tickVisual}/\texttt{snapVisualToTargets}; rectangular target rebuild
uses one \texttt{includeRect}. \texttt{GLTexture::UpdateSubImage} packs partial
dirty rects tightly and maps only the staging range (no full-buffer orphan
every frame).\\
\textbf{Why:} Per-cell div/mod and full-texture PBO orphan were measurable on
large dirty rects while fade quality must stay.\\
\textbf{Consequences:} Fade math unchanged; MockBackend token path unaffected.
\end{decisionbox}

%----------------------------------------------------------------------------
\subsection*{D-P7 --- Optional row-parallel generation}
\begin{decisionbox}[title={D-P7 \quad 2026-08-06}]
\textbf{Decision:} \texttt{RuleSet::calcGeneration} may partition rows across
workers when cell count $\ge 512\times 512$ (auto) or when
\texttt{setWorkerOverride(N)} forces $N$ workers. Results are bit-identical to
serial. Presentation/fade stays main-thread.\\
\textbf{Why:} Large dense grids are eval-bound after D-P5; spawn cost makes
parallel unhelpful on default $80\times 60$.\\
\textbf{Consequences:} Tests cover serial/parallel identity
(\texttt{Illumo.Sim.SerialParallelIdentical}); micro-benches under
\texttt{Illumo.Sim.MicroBench}. Persistent thread pool deferred.
\end{decisionbox}

%----------------------------------------------------------------------------
\subsection*{D-P4 --- dirty-rect PBO uploads (R8 path partially superseded)}
\begin{decisionbox}[title={D-P4 \quad 2026-07-28}]
\textbf{Decision (original):} Present cells as \texttt{GL\_R8} texture of state
bytes plus a $256\times 1$ RGB palette (\texttt{evalCell}). Upload dirty rects
via PBO ping-pong in \texttt{GLTexture::UpdateSubImage}. Drop dual float RGB
staging.\\
\textbf{Why:} Large grids: bandwidth and CPU of RGB float fade dominated.\\
\textbf{Consequences (kept):} dirty-rect \texttt{UpdateTexture}, PBO path,
bind-state tracker in \symbolref{GLDevice}.\\
\textbf{Superseded presentation (2026-07-30):} live path restored to CPU RGB
fade (\texttt{targetRgb}/\texttt{displayRgb}) + RGB display texture so cell
colors ease rather than snap. R8+palette is historical experiment, not current
code truth. Canonical: \pathref{docs/architecture-consensus.md}.
\end{decisionbox}

%----------------------------------------------------------------------------
\subsection*{D-E1 --- Module registration lives in App}
\begin{decisionbox}[title={D-E1 \quad 2026-07-28}]
\textbf{Decision:} \symbolref{Illumo} does not include or construct
\symbolref{CellGameModule}. \pathref{App/CellMain.cpp} calls \texttt{Init},
\texttt{addModule}, then \texttt{StartModules}. Debug module is registered
from App under \texttt{\#ifndef NDEBUG}.\\
\textbf{Why:} Engine must not depend on Game headers; product composition
belongs at the application boundary.\\
\textbf{Consequences:} New products/tests can register different module sets;
Engine only knows \symbolref{IModule}.
\end{decisionbox}

%----------------------------------------------------------------------------
\subsection*{D-E2 --- InputManager has no Game types}
\begin{decisionbox}[title={D-E2 \quad 2026-07-28}]
\textbf{Decision:} Remove \texttt{CellContext*} and Game includes from
\symbolref{InputManager}. Also drop unused \texttt{SplashText} extern from the
input header.\\
\textbf{Why:} Services must not depend on Game; input is a pure device/action
service. Splash and rules stay in Game/Rendering.\\
\textbf{Consequences:} InputManager header only needs KeyCode/InputContext/GLFW.
\end{decisionbox}

%----------------------------------------------------------------------------
\subsection*{D-E3 --- EntityTable is not the cell path}
\begin{decisionbox}[title={D-E3 \quad 2026-07-28}]
\textbf{Decision:} Archive \pathref{EntityTable.h} and \pathref{ModuleObject.h}
under \pathref{archive/dead-engine/}. Remove \texttt{entityTable} from
\symbolref{Scene}/\symbolref{IllumoContext}.\\
\textbf{Why:} Nothing allocated or used an EntityTable; the live path is
drawables + tokens, not ECS.\\
\textbf{Consequences:} Reintroduce a general object table only when a real
feature needs it; document that cells are not entities.
\end{decisionbox}

%----------------------------------------------------------------------------
\subsection*{D-E4 --- Scene is a drawable list, not a graph}
\begin{decisionbox}[title={D-E4 \quad 2026-07-28}]
\textbf{Decision:} \symbolref{Scene} holds only \texttt{drawables},
\texttt{window}, and \texttt{activeCamera}. Remove \texttt{root},
\texttt{nodeLookup}, no-op \texttt{Update}, and unused
\texttt{RenderableObject}. Archive \pathref{SceneObject.h}; \symbolref{Camera}
and \symbolref{CommandLine} no longer inherit a node graph.\\
\textbf{Why:} Architecture review (dead graph scaffolding); production never
used parent/child transforms for cells or UI.\\
\textbf{Consequences:} Frame contribution model is explicit and small.
\end{decisionbox}

%----------------------------------------------------------------------------
\subsection*{D-E5 --- Freeze IllumoContext; validate at Start}
\begin{decisionbox}[title={D-E5 \quad 2026-07-28}]
\textbf{Decision:} Keep \symbolref{IllumoContext} as a service bag for the
\textbf{two shipped modules} only. Document required fields per module.
\texttt{Start()} logs an error and bails if the bag is incomplete.
Do not grow the bag for convenience; a third module should take explicit
constructor dependencies if its subset differs.\\
\textbf{Why:} Architecture review --- bag cost is zero at $N=2$ but becomes a
de-facto global API as $N$ grows.\\
\textbf{Consequences:} \texttt{IllumoContextHasGameCore} /
\texttt{IllumoContextHasDebugCore} helpers; freeze note on the struct.
\end{decisionbox}

%----------------------------------------------------------------------------
\subsection*{D-C1 --- Canvas dual role is intentional (for now)}
\begin{decisionbox}[title={D-C1 \quad 2026-07-28}]
\textbf{Decision:} Keep \symbolref{Canvas} as domain grid + fade view + GPU
enroll in one type. Document the scale wall: full-grid
\texttt{calcGeneration}, dual float RGB buffers ($\sim$24\,B/cell), no
chunked simulation domain.\\
\textbf{Why:} Locality for dirty-rect + palette + fade; headless tests already
use injected MockBackend rather than pure LifeGrid.\\
\textbf{Consequences:} Split \texttt{LifeGrid}/\texttt{CanvasView} only if pure
sim tests or large canvases become product requirements.
\end{decisionbox}

%----------------------------------------------------------------------------
\subsection*{D-R9 --- String-named uniforms are backend debt}
\begin{decisionbox}[title={D-R9 \quad 2026-07-28}]
\textbf{Decision:} Keep \texttt{char name[32]} uniforms for OpenGL + Mock.
Treat as debt if a second \emph{real} GPU backend is adopted; prefer binding
points / pre-resolved locations then.\\
\textbf{Why:} Architecture review --- uniform path is the least portable part
of \symbolref{RenderCommand}; handles and UpdateTexture are closer to neutral.\\
\textbf{Consequences:} Comment on uniform payloads; no API change now.
\end{decisionbox}

%----------------------------------------------------------------------------
\subsection*{D-R10 --- Production drawables are pure-token}
\begin{decisionbox}[title={D-R10 \quad 2026-07-28}]
\textbf{Decision:} \symbolref{Canvas}, \symbolref{CommandLine},
\symbolref{GLString}, \symbolref{SplashText} use
\texttt{AppendCommands} only (return true when contributing). Immediate
\texttt{Draw()} hybrid remains for tests/stubs. CRTP \texttt{DrawImpl} is
legacy for that fallback.\\
\textbf{Why:} Migration complete for shipped UI/canvas; dual path is debt only
where unmigrated types exist.\\
\textbf{Consequences:} Documented in \pathref{Drawable.h} and
\texttt{Renderer::RenderScene}.
\end{decisionbox}

%----------------------------------------------------------------------------
\subsection*{D-R14 --- Scene layers and built-in render styles}
\begin{decisionbox}[title={D-R14 \quad 2026-08-07}]
\textbf{Decision:} Introduce \textbf{Phase A} composition without multi-pass
GPU targets:
\begin{itemize}
  \item \symbolref{Scene} buckets drawables into ordered layers
        (\texttt{World}, \texttt{UI}, \texttt{Debug}) within one main pass
        (default framebuffer, single clear).
  \item \symbolref{Renderer} owns built-in \symbolref{RenderStyle} entries
        (\texttt{Canvas}, \texttt{UiText}, \texttt{Console}): each holds a
        shader table handle plus \texttt{PipelineState} defaults.
  \item Production drawables keep content handles (mesh/texture/domain data)
        and call \texttt{bindStyle} before content tokens.
\end{itemize}
Layers are not the same as \symbolref{RenderPass} (targets/load-store). A
layer may later feed multiple passes if targets diverge; that is deferred.\\
\textbf{Why:} World/UI/debug stacking and shared draw policy should not live
as duplicated per-drawable shader enrollment and ad-hoc pipeline structs.\\
\textbf{Consequences:} \pathref{RenderLayerId.h}, \pathref{RenderStyle.h},
\pathref{RendererStyles.cpp}; modules pass a layer to
\texttt{Scene::AddDrawable}; docs distinguish layer vs pass.
\end{decisionbox}

%----------------------------------------------------------------------------
\subsection*{D-R15 --- Render primitives and GameVisual host}
\begin{decisionbox}[title={D-R15 \quad 2026-08-07}]
\textbf{Decision:} Introduce value-type \textbf{render primitives}
(\symbolref{ShapePrimitive}, \symbolref{SpritePrimitive}) composed on a
\symbolref{GameVisual} \texttt{Drawable} host.
\begin{itemize}
  \item Primitives are not Scene entries; the host is.
  \item Game/editor objects \textbf{embed} \symbolref{GameVisual} (or multiple)
        and add primitives to build complex visuals without hand-rolling
        enroll/bind/draw tokens.
  \item Built-in styles \texttt{Shape} and \texttt{Sprite} live on
        \symbolref{Renderer}; GPU mesh/texture/shader objects stay in the
        backend registries.
  \item Coordinate space is host policy (\texttt{Pixels} or \texttt{World}).
\end{itemize}
\textbf{Why:} Composition over per-primitive drawables keeps Scene traffic low
and matches D-R14 style ownership.\\
\textbf{Consequences:} \pathref{Rendering/Primitives/}; headless tests
\texttt{Illumo.GameVisual.*}; Canvas/UI paths unchanged.
\end{decisionbox}

%----------------------------------------------------------------------------
\subsection*{D-F1 --- MacroDefs Windows include deferred}
\begin{decisionbox}[title={D-F1 \quad 2026-07-30}]
\textbf{Decision:} Do \textbf{not} split \pathref{MacroDefs.h} /
\texttt{Windows.h} now. Revisit only if include toxicity causes real pain
(compile time, macro collisions, portability friction).\\
\textbf{Why:} Author priority --- convenience of \texttt{\_\_SLEEP} outweighs
hygiene until symptoms appear.\\
\textbf{Consequences:} Status note in services/platform section; not an open
blocker.
\end{decisionbox}

%----------------------------------------------------------------------------
\subsection*{D-N1 --- Illumo is the project and product name}
\begin{decisionbox}[title={D-N1 \quad 2026-08-04}]
\textbf{Decision:} Adopt \textbf{Illumo} as the repository-facing project,
product, CMake target, runtime title, console name, and save-format name. Keep
the existing \symbolref{Illumo} runtime-host class; the shared name is
intentional. Canonical tests are \texttt{IllumoTests}, and new saves use the
\texttt{.illumo} extension.\\
\textbf{Why:} The author selected Illumo as the permanent product identity.\\
\textbf{Consequences:} Live first-party source and documentation use Illumo;
generated build trees, vendored material, and archived historical artifacts are
not implementation sources and are not rewritten by this decision.
\end{decisionbox}

%----------------------------------------------------------------------------
\subsection*{D-B1 --- DebugModule is excluded from Release}
\begin{decisionbox}[title={D-B1 \quad 2026-08-04}]
\textbf{Decision:} \symbolref{DebugModule} is Debug-build product composition
only. \symbolref{CellMain} must include/register it behind
\texttt{\#ifndef NDEBUG}, and CMake must compile its implementation only for the
Debug configuration.\\
\textbf{Why:} The developer overlay, console input handling, FPS/debug labels,
and instrumentation are development tools, not Release product behavior.\\
\textbf{Consequences:} Shared command/logging service code may remain available
to tests or non-UI systems, but Release must contain no DebugModule object or
module registration. Audit this boundary after CMake/composition changes.
\end{decisionbox}

%----------------------------------------------------------------------------
\subsection*{D-CLI1 --- Split generic and domain console ownership}
\begin{decisionbox}[title={D-CLI1 \quad 2026-08-04}]
\textbf{Decision:} \symbolref{CommandLine} owns editing, parsing, help,
environment/application commands, and rendering. Product modules register
domain behavior through \symbolref{CommandRegistry}; each registration carries
usage, description, and optional completion candidates.\\
\textbf{Why:} Services should not acquire Game dependencies, while help and
completion must describe the commands that actually exist.\\
\textbf{Consequences:} \symbolref{CellGameModule} owns simulation, canvas,
ruleset, camera, and save/load commands. The shared known-ruleset list drives
validation/completion. Queued commands return successfully without falling
through to environment lookup or an unknown-command error.
\end{decisionbox}

%----------------------------------------------------------------------------
\subsection*{D-UI1 --- Measured console editor in one bounded batch}
\begin{decisionbox}[title={D-UI1 \quad 2026-08-04}]
\textbf{Decision:} Render the command caret and selection from measured text
prefixes, keep long input visible through a horizontal viewport, and retain one
batched token draw for console chrome and text. Size the dynamic mesh for 6,000
UI quads and regression-test a detailed help page beyond the former cap.\\
\textbf{Why:} An appended underscore produced an apparent cursor offset, while
the 2,000-quad mesh silently truncated easy-font output because each character
uses several quads.\\
\textbf{Consequences:} Editing supports selection and word operations without
lying about insertion position; a full help page remains one draw without
cutting off its final line.
\end{decisionbox}

%----------------------------------------------------------------------------
\subsection*{D-UI2 --- Wrapped console history with unified scroll metrics}
\begin{decisionbox}[title={D-UI2 \quad 2026-08-06}]
\textbf{Decision:} Word-wrap console history to the panel width, scroll by
visual lines rather than raw history entries, and compute PageUp/wheel/scrollbar
limits from the same floating-aware panel layout used when drawing. Raise the
console UI mesh to 8,000 quads to leave headroom for wrapped lines plus chrome.\\
\textbf{Why:} Hard truncation hid the rest of long help lines, and scroll
handlers used a mounted-only height estimate that disagreed with the floating
panel, so paging to the oldest history could leave text clipped or unreachable.\\
\textbf{Consequences:} Scrolling fully to the start of history shows the oldest
characters completely; floating resize and mounted mode share one max-scroll
definition; a focused wrap/scroll regression covers the clamp and batch path.
\end{decisionbox}

%----------------------------------------------------------------------------
\subsection*{D-DOC1 --- One first-party documentation tree}
\begin{decisionbox}[title={D-DOC1 \quad 2026-08-04}]
\textbf{Decision:} Keep all first-party project documentation under
\pathref{docs/}: current Markdown at the top level, package maps in
\pathref{packages/}, dated records in \pathref{sessions/}, provenance in
\pathref{history/}, one current LaTeX entrypoint at \pathref{latex/illumo.tex},
and generated output in \pathref{output/}.\\
\textbf{Why:} Competing LaTeX entrypoints and package READMEs scattered beside
code made current truth and PDF generation difficult to discover.\\
\textbf{Consequences:} The repository root keeps only the landing README and
AGENTS bootstrap as intentional prose exceptions. Licenses and vendored notes
remain beside governed assets. Architecture changes update Markdown, relevant
LaTeX chapters, and this append-only log together. \pathref{docs/build.ps1}
is the PDF wrapper and the optional default-build \texttt{IllumoDocs} target
uses that same script when TeX is available.
\end{decisionbox}

%----------------------------------------------------------------------------
\subsection*{D-T1 --- Granular CTest registration and coverage gate}
\begin{decisionbox}[title={D-T1 \quad 2026-08-04}]
\textbf{Decision:} Keep one \texttt{IllumoTests} executable for compile
efficiency, but register each logical behavior as an exact
\texttt{Illumo.<area>.<case>} CTest entry. CTest runs one filtered case per
process in an isolated working directory. With
\texttt{ILLUMO\_ENABLE\_COVERAGE=ON}, Clang/LLVM adds the
\texttt{IllumoCoverage} target and enforces at least 85\% line coverage over
headless-testable first-party production code.\\
\textbf{Why:} A monolithic CTest result hid which of dozens of behaviors ran or
failed, and a narrow denominator could report a strong percentage while omitting
game commands, persistence, input, or assets.\\
\textbf{Consequences:} New behaviors receive separate stable names; file-backed
tests may run in parallel without sharing working files; tests, vendored/system
code, MockBackend, and the live OpenGL backend are excluded from the automated
production denominator. Native dialogs, real-window behavior, and live OpenGL
still require proportional smoke tests. This supersedes the single-CTest-name
registration detail in D-R7 without changing the MockBackend decision.
\end{decisionbox}

%----------------------------------------------------------------------------
\subsection*{D-UI3 --- Floating mode, title-bar dragging, and corner resize}
\begin{decisionbox}[title={D-UI3 \quad 2026-08-06}]
\textbf{Decision:} Extend \symbolref{CommandLine} to support dynamic window modes and resizable floating window geometry. Mode switching is driven by double-clicking the console title bar or executing \texttt{console\_mode [floating|mounted|toggle]}. In floating mode, left-clicking and dragging the title bar repositions the window across the screen, and dragging the bottom-right corner grip handle (or running \texttt{console\_size <W> <H> | reset}) resizes the console window dynamically with real-time viewport clipping and single-batch quad rendering.\\
\textbf{Why:} Enhances Developer Console usability, allows moving and scaling the console UI away from active grid regions, and provides modular window layout flexibility.\\
\textbf{Consequences:} \symbolref{CommandLine} geometry (chrome, history text, input trough, scrollbar, ghost text, selection highlight) is computed relative to dynamic floating bounds \texttt{panelX0}, \texttt{panelX1}, \texttt{panelY0}, \texttt{panelY1}; mouse press/drag event dispatch accounts for floating position offset and corner resize bounds.
\end{decisionbox}

%----------------------------------------------------------------------------
\subsection*{D-R11 --- Backend constructed at composition root}
\begin{decisionbox}[title={D-R11 \quad 2026-08-06}]
\textbf{Decision:} Construct the concrete production backend
(\symbolref{GLBackend}) in \symbolref{Illumo::Init} and inject it into
\symbolref{Renderer} as \texttt{std::unique\_ptr<IBackend>}. \symbolref{Renderer}
must not include or construct OpenGL types. Tests continue to inject
\symbolref{MockBackend} through the same ownership path (or a non-owning raw
pointer for stack fixtures).\\
\textbf{Why:} Dependency inversion: \texttt{Renderer $\rightarrow$ IBackend},
composition root $\rightarrow$ concrete backend. Production and headless test
architectures become structurally identical.\\
\textbf{Consequences:} \pathref{Renderer.h} no longer includes
\pathref{OpenGL/GLBackend.h}; OpenGL coupling stays under
\pathref{Rendering/OpenGL/} and the host composition file.
\end{decisionbox}

%----------------------------------------------------------------------------
\subsection*{D-R12 --- CommandQueue overflow policy}
\begin{decisionbox}[title={D-R12 \quad 2026-08-06}]
\textbf{Decision:} Keep the fixed 2048-token \symbolref{CommandQueue} capacity.
When full, drop additional tokens, count drops, and log an error once per
frame (once per \texttt{Reset} interval) rather than failing silently or
growing the buffer unbounded.\\
\textbf{Why:} Silent loss hid runaway emitters; unbounded growth is not
acceptable on the production frame path. One log per overflowing frame is
visible without spamming every discarded command.\\
\textbf{Consequences:} Callers can query drop counters in tests; capacity
changes require profiling evidence, not opportunistic growth.
\end{decisionbox}

%----------------------------------------------------------------------------
\subsection*{D-WW1 --- Wireworld seed and edit brush}
\begin{decisionbox}[title={D-WW1 \quad 2026-08-06}]
\textbf{Decision:} Startup seeding is ruleset-aware. Binary life-like modes seed
a centered Game-of-Life glider. Wireworld seeds a horizontal conductor line
with a head+tail electron pair. In Wireworld edit mode, left-paint uses a
sticky brush selected by keys \texttt{1}/\texttt{H} (head), \texttt{2} (empty),
\texttt{3}/\texttt{T} (tail), \texttt{4} (conductor); right-click paints empty.\\
\textbf{Why:} The previous GoL-only glider produced five heads under Wireworld,
and mouse paint could not place heads without the console \texttt{setcell}
command.\\
\textbf{Consequences:} Wireworld is usable for circuit editing without leaving
edit mode; focused tests cover seed cells and brush selection state.
\end{decisionbox}

%----------------------------------------------------------------------------
\subsection*{D-C2 --- CellGrid domain extracted from Canvas}
\begin{decisionbox}[title={D-C2 \quad 2026-08-06}]
\textbf{Decision:} Introduce \symbolref{CellGrid} as pure dense simulation
storage (life buffer + dirty tracking) with no \symbolref{Renderer}, window, or
camera dependency. \symbolref{Canvas} extends \symbolref{CellGrid} for
presentation (RGB fade, texture upload) and GPU enrollment.
\symbolref{RuleSet} operates on \texttt{CellGrid*} only.\\
\textbf{Why:} The architecture assessment required a truthful domain boundary so
simulation can be constructed, stepped, and read without a graphics stack.
D-C1's intentional monolith is refined, not discarded: presentation still lives
on \symbolref{Canvas}.\\
\textbf{Consequences:} Headless domain tests construct \symbolref{CellGrid}
(or a domain-only \symbolref{Canvas} with null render services) and drive
shipped \symbolref{RuleSet} implementations. A further
\texttt{CanvasRenderer} split remains optional.\\
\textbf{Supersedes / refines:} D-C1 (monolith until pain).
\end{decisionbox}

%----------------------------------------------------------------------------
\subsection*{D-R13 --- Single production frame path}
\begin{decisionbox}[title={D-R13 \quad 2026-08-06}]
\textbf{Decision:} \symbolref{Illumo::Render} always rebuilds the per-frame
drawable list and submits tokens through \symbolref{Renderer::RenderScene}.
Remove the env \texttt{UseTokenProof} product bypass. Keep
\symbolref{Renderer::RenderProofQuad} for headless token e2e tests only.
Hybrid immediate \texttt{Draw()} remains only as a stub/test fallback when
\texttt{AppendCommands} returns false.\\
\textbf{Why:} Competing production frame paths create ambiguous ordering and
test gaps.\\
\textbf{Consequences:} One default product path; proof coverage stays in
\texttt{IllumoTests} without affecting the live loop.
\end{decisionbox}

%----------------------------------------------------------------------------
\subsection*{D-OWN1 --- Module-owned mode splash}
\begin{decisionbox}[title={D-OWN1 \quad 2026-08-06}]
\textbf{Decision:} Own the EDIT/NORMAL splash label as
\texttt{std::unique\_ptr<SplashText> modeSplash} on
\symbolref{CellGameModule}. Delete the file-scope \texttt{stateSplash}
mutable global.\\
\textbf{Why:} File-scope mutable product UI hid ownership and complicated
teardown and multi-instance reasoning.\\
\textbf{Consequences:} Splash lifetime follows the game module; platform
\texttt{SaveLoad} dialog entry points remain platform implementations without
becoming a new engine singleton.
\end{decisionbox}

%----------------------------------------------------------------------------
\subsection*{Template for new entries}
\begin{lstlisting}[language={}]
\subsection*{D-0XX --- short title}
\begin{decisionbox}[title={D-0XX \quad YYYY-MM-DD}]
\textbf{Decision:} ...
\textbf{Why:} ...
\textbf{Consequences:} ...
\textbf{Supersedes:} D-0YY (optional)
\end{decisionbox}
\end{lstlisting}
